\section{Trabalhos Relacionados}
\label{sec:trabalhos_relacionados}

Esta seção apresenta estudos sobre a qualidade das descrições MCP, o desempenho de agentes no ecossistema e o uso de LLMs como avaliadores, que fundamentam a lacuna desta pesquisa.

Hasan et al. (2026) \nocite{2026arXiv260214878M} investigaram a qualidade das documentações de 856 ferramentas no ecossistema MCP por meio de uma análise empírica fundamentada em LLMs. O estudo aplicou uma rubrica de seis dimensões focada na clareza e completude do conteúdo textual. Os resultados revelaram que 97,1\% das descrições possuem ao menos um \textit{smell}, o que reduz a precisão dos agentes na seleção de parâmetros. \textcolor{red}{O presente estudo parte dessa investigação e avança em uma direção não explorada pelos autores: a incorporação do código ao processo de avaliação, verificando em que medida o acesso à lógica interna do servidor altera os escores obtidos nas dimensões da rubrica proposta pelos autores}.

Shi et al. (2026) \nocite{shi2026description} investigam a prevalência e os impactos de inconsistências semânticas entre descrições de ferramentas MCP e seus respectivos códigos, propondo uma taxonomia de inconsistências e um detector baseado em LLM aplicado em larga escala a servidores reais. Os resultados revelaram que 9,93\% das ferramentas analisadas apresentam inconsistências e que 35\% dos servidores possuem ao menos um caso, com riscos de seleção incorreta de ferramentas, efeitos colaterais não declarados e falhas de segurança. \textcolor{red}{Esse estudo sustenta a premissa desta pesquisa ao demonstrar que informações relevantes sobre o comportamento de uma ferramenta frequentemente não estão presentes em sua descrição.}

Guo et al. (2025)\nocite{guo2025measurement} mapeiam a adoção do MCP com um rastreador que analisa mais de 8.400 projetos e constatam que mais da metade dos projetos listados são inválidos, abandonados ou de baixo valor agregado. O estudo caracteriza os repositórios, mas não examina a consistência entre descrições e implementação; o dado de baixa maturidade também fundamenta o critério de popularidade adotado na seleção da amostra (Seção~\ref{sec:materiais}).

Luo et al. (2025)\nocite{2025arXiv250814704L} propuseram o MCP-Universe, um \textit{benchmark} que avalia modelos de linguagem em servidores MCP reais por meio de avaliadores baseados em execução. Ao avaliar mais de 20 modelos em 231 tarefas distribuídas por 11 servidores, os resultados revelaram que mesmo modelos de fronteira atingem taxas de sucesso abaixo de 45\%, com falhas frequentemente associadas ao uso incorreto de interfaces de ferramentas pouco familiares. \textcolor{orange}{O presente trabalho não mede o desempenho de agentes; ele examina um fator anterior a essas falhas, a qualidade da informação que a descrição oferece ao agente sobre a interface da ferramenta, e se o código revela omissões nessa informação.}

Zheng et al. (2023) \nocite{NEURIPS2023_91f18a12} investigam a viabilidade da abordagem \textit{LLM-as-a-Judge}, comparando julgamentos do GPT-4 com avaliações humanas em \textit{benchmarks} de conversação, e mostram concordância superior a 80\%, comparável à observada entre os próprios especialistas. Esse resultado fundamenta o uso de avaliadores baseados em LLM neste estudo.

Em síntese, \textcolor{orange}{é} demonstrado que a qualidade das descrições MCP é um fator crítico e que texto e implementação podem divergir, mas as abordagens de avaliação de qualidade analisam as documentações de forma isolada. Este trabalho investiga como a incorporação do código influencia essa avaliação.
